\section{General robust repair}

The profile envelope and parameterized exchange theorem are the general
repair statements of the paper.  We first record a secondary saturation
consequence of the carrier criterion, then derive those two results.

\subsection*{A saturation consequence}

Call an invariant arc \emph{pair-saturated} if it has no legal conjugate-pair
extension.  This is weaker than ordinary completeness.

\begin{corollary}[Orbit-saturation bound]\label{cor:saturation}
If a $\phi$-invariant $k$-arc in $\PG(2,s^2)$ is pair-saturated, then
\begin{equation}\label{eq:saturation}
  k\ge 1+\left\lceil\sqrt{2s(s-1)}\right\rceil.
\end{equation}
In particular, for every prime power $s\ge7$, every invariant eight-arc in
$\PG(2,s^2)$ admits a conjugate-pair extension.
\end{corollary}

\begin{proof}
If every fixed line is occupied, Lemma~\ref{lem:empty-lines} gives
$s^2+s+1=f(s+1)-\binom f2+e$.  Substituting $2e=k-f$ and completing the
square gives the occupation identity
\[
  k=s^2+1+(f-s-1)^2,
\]
which is stronger than \eqref{eq:saturation}.  Otherwise choose an empty
fixed line.  Pair saturation forces all $N=s(s-1)/2$ of its candidates to
be forbidden, so $N\le M$.  From $k=f+2e$ and \eqref{eq:M},
\[
  M=e((k-1)-e),
  \qquad
  4M\le(k-1)^2.
\]
Consequently $2s(s-1)\le(k-1)^2$, which gives
\eqref{eq:saturation}.

For $k=8$, one has $M\le12<N$ when $s\ge7$.  Full occupation is impossible
because $8<s^2+1$, so Theorem~\ref{thm:pair-criterion} applies.
\end{proof}

\subsection*{Profile envelope and robust exchange}

\begin{theorem}[Five-profile lower envelope]\label{thm:profile-envelope}
Let $s\ge7$ and let $D$ be a $\phi$-invariant eight-arc with profile
$8=f+2e$.  Put $N=s(s-1)/2$.  Then $\Npair(D)\ge L_f(s)$, where
\[
\begin{array}{c|c|c|c}
(f,e)&E_f(s)&M_f&L_f(s)=E_f(s)(N-M_f)\\ \hline
(0,4)&s^2+s-3&12&(s^2+s-3)(N-12)\\
(2,3)&s^2-s-3&12&(s^2-s-3)(N-12)\\
(4,2)&s^2-3s+1&10&(s^2-3s+1)(N-10)\\
(6,1)&s^2-5s+9&6&(s^2-5s+9)(N-6)\\
(8,0)&s^2-7s+21&0&(s^2-7s+21)N.
\end{array}
\]
For the certified first-order envelope $L(s)=\min_f L_f(s)$,
\[
  L(7)=319,\qquad L(8)=726,\qquad L(9)=1350,
\]
and $L(s)=L_8(s)$ for every $s\ge10$.  The minimizing profiles are therefore
$f=4$ at $s=7$, $f=6$ at $s=8,9$, and $f=8$ for $s\ge10$.
In particular, $\Npair(D)\ge319$ for every $s\ge7$.
\end{theorem}

\begin{proof}
Substituting the five solutions of $8=f+2e$ into
\eqref{eq:E} and \eqref{eq:M} gives the displayed values of
$E_f$, $M_f$, and hence $L_f=E_f(N-M_f)$ by
Theorem~\ref{thm:pair-criterion}.  At $s=7,8,9$ the five value vectors are
respectively
\[
(477,351,319,345,441),\quad
(1104,848,738,726,812),\quad
(2088,1656,1430,1350,1404).
\]
For $s=t+10$, twice each difference $L_f-L_8$, for
$f=0,2,4,6$, is respectively
\[
\begin{aligned}
 8t^3+184t^2+1280t+2472,&\qquad
 6t^3+126t^2+768t+1152,\\
 4t^3+76t^2+400t+380,&\qquad
 2t^3+34t^2+152t+12.
\end{aligned}
\]
All coefficients are nonnegative for $t\ge0$.
This proves the infinite branch.  Independently, each $L_f(s)$ is
nondecreasing for $s\ge7$, so the value $319$ at the first boundary is a
uniform lower bound.  These polynomial identities and inequalities are
kernel-checked in the supplementary Lean source.
\end{proof}

\begin{corollary}[Quantitative alternate-orbit repair]
\label{cor:alternate-repair}
Let $s\ge7$ be a prime power.  If $A$ is a $\phi$-invariant ten-arc in
$\PG(2,s^2)$ and $O\subseteq A$ is any selected nonfixed conjugate orbit,
then deleting $O$ leaves at least $318$ alternate-orbit repairs.
\end{corollary}

\begin{proof}
Put $D=A\setminus O$.  Then $D$ is an invariant eight-arc, and $O$ is one
of its legal conjugate-pair extensions.  Theorem~\ref{thm:profile-envelope}
gives $\Npair(D)\ge319$.  Excluding $O$ leaves at least $318$ different
legal pairs.
\end{proof}

\begin{theorem}[Parameterized robust equivariant exchange]
\label{thm:parameterized-repair}
Let $s\ge3$ be a prime power, let $k\ge1$, and let $r\ge0$.  Put
\[
  N(s)=\frac{s(s-1)}2,
  \qquad
  W(k)=\left\lfloor\frac{(k-1)^2}{4}\right\rfloor.
\]
If
\begin{equation}\label{eq:repair-phase}
  W(k)+r+1\le N(s),
\end{equation}
then deleting any selected nonfixed conjugate orbit from a $\phi$-invariant
$(k+2)$-arc in $\PG(2,s^2)$ leaves at least $r$ alternate-orbit repairs.
The obstruction envelope $W(k)$ is the exact arithmetic maximum over all
profiles $k=f+2e$; no geometric realization of every maximizing profile is
asserted.
In particular, if $s\ge4$ and $1\le k\le s+1$, every such deletion has an
alternate repair.
\end{theorem}

\begin{proof}
For a remainder profile $k=f+2e$, the noninvariant old-secant obstruction is
\[
  M=f e+e(e-1)=e(k-e-1).
\]
The concave quadratic has the exact profile maximum
$M\le W(k)$: equality is attained by $f=0$ when $k$ is even and by $f=1$
when $k$ is odd.  Condition~\eqref{eq:repair-phase} implies $k\le2s+1$.
For $s\ge3$ this gives $k<s^2+1$; the completed-square occupation identity
\[
  2E_{\mathrm{occ}}=(s+1)^2+k-(f-s-1)^2
\]
then shows that at least one fixed carrier is empty.  On that carrier,
Theorem~\ref{thm:pair-criterion} gives
\[
  \Npair(D)\ge N(s)-M\ge N(s)-W(k)\ge r+1.
\]
The erased orbit is one of these legal pairs, so excluding it leaves at
least $r$ alternatives.

For the rectangular corollary, $k\le s+1$ gives
$4W(k)\le(k-1)^2\le s^2$, while $s\ge4$ gives
$s^2+8\le2(s^2-s)=4N(s)$.  Hence $W(k)+2\le N(s)$, which is
\eqref{eq:repair-phase} with $r=1$.  The restriction $s\ge4$ is genuine for
this rectangle: $(s,k)=(3,4)$ misses the inequality by one.
\end{proof}

\begin{corollary}[All base orders from five]\label{cor:all-orders}
Let $s\ge5$ be a prime power.  Every Frobenius-invariant eight-arc in
$\PG(2,s^2)$ admits an extension by a fresh nonfixed point together with
its Frobenius conjugate.
\end{corollary}

\begin{proof}
The case $s=5$ is Theorem~\ref{thm:q25}.  There is no prime power strictly
between $5$ and $7$, and the final assertion of
Corollary~\ref{cor:saturation} treats every $s\ge7$.
\end{proof}

\begin{proof}[Proof of Theorem~\ref{thm:main}]
The extension assertion is Corollary~\ref{cor:all-orders}.
Theorem~\ref{thm:q25} and Corollary~\ref{cor:q25-uniform-repair} give (i),
including the exceptional profile supplied by Proposition~\ref{prop:f2}.
Theorem~\ref{thm:profile-envelope} and
Corollary~\ref{cor:alternate-repair} give (ii).
\end{proof}

The scale in \eqref{eq:saturation} is the classical Lunelli--Sce scale for
ordinary complete arcs, with the same secant-covering mechanism in the
Frobenius quotient.  Ng and Wild's line-covering problem
\cite{NgWild2001} is another adjacent square-root-scale result.  The point
of Corollary~\ref{cor:saturation} is the weaker pair-saturation hypothesis,
not a new constant or a sharp family.
